\subsection{Choice, Zorn's Lemma, and Well-Ordering}

\begin{principle} \label{princ:axiom-of-choice-analysis} \hlt{Choice Principle}\\
For nonempty sets $\{X_i\}_{i\in I}$, the axiom of choice asserts $\exists c:I\to\bigcup\limits_{i\in I}X_i$ such that $c(i)\in X_i\ \forall i\in I$.\\
Use direct construction for finite or explicit choices; use choice only for arbitrary simultaneous choices.
\end{principle}

\begin{principle} \label{princ:zorns-lemma} \hlt{Zorn's Lemma}\\
If a nonempty partially ordered set $X$ has the property that every chain in $X$ has an upper bound in $X$, then $X$ contains a maximal element. This is equivalent in strength to the axiom of choice, Principle \ref{princ:axiom-of-choice-analysis}.
\end{principle}

\begin{method} \label{meth:cardinality-proof-patterns}\label{meth:cardinality-proof-patterns-source-todo} \hlt{Cardinality Proof Patterns}
\begin{enumerate}[label=(\roman*)]
\setlength{\itemsep}{0pt}
\item Same cardinality: construct a bijection.
\item Countable: construct a sequence whose range contains every element.
\item Uncountable: show every proposed sequence misses some element.
\item Choice/maximality: cite Principle \ref{princ:axiom-of-choice-analysis} for arbitrary simultaneous choices and Principle \ref{princ:zorns-lemma} for maximal elements in partially ordered sets.
\end{enumerate}
\hlt{Key Idea}: Use maps: bijections prove equal cardinality, enumerations $\N\to A$ prove countability, and diagonal arguments show every proposed enumeration misses an element.
\end{method}

\hlt{Placeholder}: Add the well-ordering theorem, maximality methods, and examples of where choice is avoidable versus structurally necessary.
